\section{Introduction}
Distinguishing salivary gland mucosa-associated lymphoid tissue (MALT) lymphoma from lymphoid-rich inflammation associated with Sjögren's syndrome (SS) requires assessment of the infiltrate in its glandular context. SS is an autoimmune setting in which MALT lymphoma can arise, and lymphoepithelial lesions or architectural disturbance alone do not resolve every differential diagnosis \citep{ref1,ref2}. Integrated assessment may require morphology, immunophenotyping and molecular findings within the relevant lymphoma classification framework \citep{ref3,ref4}. A computational analysis of routine whole-slide images could provide a quantitative view of this material. The immediate research question is whether slide images contain a reproducible classification signal for the two recorded diagnostic categories.

Disease-focused digital pathology has begun to address related questions. Cell--tissue graph models have been used for SS classification, and an attention model evaluated whether minor labial salivary gland biopsies had a focus score of at least one \citep{ref10,focus2025}. Other studies classified orbital MALT lymphoma versus IgG4-related ophthalmic disease, or distinguished MALT and diffuse large B-cell lymphoma categories using pathologist-selected tissue images \citep{ref11,lymphoma2025}. These studies establish relevant approaches, but their organs, reference outcomes or image-selection procedures differ. They do not directly answer how whole-slide, slide-label-supervised models behave for salivary gland MALT versus SS.

Multiple-instance learning (MIL) connects slide-level labels to collections of image patches without requiring exhaustive lesion annotations \citep{ref5,ref6,ref7}. Pathology foundation models now supply transferable patch representations, allowing downstream classifiers to be trained while keeping the image encoder fixed \citep{ref8,ref9}. This makes disease-specific experiments feasible with limited slide collections, but leaves an important empirical question: how much does a complex aggregator add beyond simpler summaries of the same features? Recent work on MIL training and feature sampling motivates separating architectural complexity from optimization and input-sampling effects \citep{nnmil2026}.

We studied retrospective MALT--SS classification using fixed UNI2-h representations and Hybrid-Evidence MIL. The analysis combines a retained model's internal validation and external performance with simpler baselines, controlled retraining and repeated patch sampling. We also examined how external archive categories related to scores and made selected image regions available for pathological review. The study evaluates slide classification and the conditions under which its measured performance changes; it does not treat a slide score as a replacement for integrated clinical diagnosis.
